\chapter{Construire et représenter une suite}
\section*{Exercice 4.1 -- Définition explicite}
Pour $u_n=(n+1)/(n+2)$, calculer $u_0,u_1,u_5$ et montrer $0<u_n<1$.

\section*{Exercice 4.2 -- Récurrence et initialisation}
Pour $u_0=1$ et $u_{n+1}=2u_n+3$, calculer jusqu'à $u_3$. Pourquoi la récurrence seule ne suffit-elle pas ?

\section*{Exercice 4.3 -- Un motif}
On aligne $n\ge1$ carrés identiques, deux voisins partageant un côté. Combien d'allumettes faut-il ? Démontrer la formule.

\section*{Exercice 4.4 -- Intervalle stable}
Pour $u_0=1$ et $u_{n+1}=\sqrt{2+u_n}$, prouver que tous les termes sont définis et appartiennent à $[1,2]$.

\section*{Exercice 4.5 -- Indices de liste}
Écrire une fonction renvoyant $u_0,\ldots,u_n$ pour $u_0=2$, $u_{k+1}=u_k^2-1$. Combien d'éléments contient la liste ?

\section*{Exercice 4.6 -- Graphique trompeur}
Les premiers termes d'une suite valent $1,2,3,4$. Peut-on affirmer $u_n=n+1$ pour tous les $n$ ? Construire un contre-exemple explicite.
